\documentclass[11pt,a4paper]{article}

\usepackage[a4paper,margin=24mm]{geometry}
\usepackage{fontspec}
\usepackage[ngerman]{babel}
\usepackage{booktabs}
\usepackage{array}
\usepackage{xcolor}
\usepackage{microtype}
\usepackage{parskip}
\usepackage{titlesec}
\usepackage{longtable}

% =========================
% Schriftarten
% =========================

\setmainfont{Inter}

\newfontfamily\headingfont{Montserrat SemiBold}

% =========================
% Farben
% =========================

\definecolor{neutralgood}{HTML}{3A6B4B}
\definecolor{neutralmid}{HTML}{8A6F2A}
\definecolor{neutraltext}{HTML}{444444}

\newcommand{\good}[1]{\textcolor{neutralgood}{#1}}
\newcommand{\neutralwarn}[1]{\textcolor{neutralmid}{#1}}
\newcommand{\neutral}[1]{\textcolor{neutraltext}{#1}}

% =========================
% Überschriften
% =========================

\titleformat{\section}
{\headingfont\fontsize{15pt}{17pt}\selectfont}
{}{0pt}{}

\pagenumbering{gobble}

% =========================
% Dokument
% =========================

\begin{document}

{\headingfont\fontsize{20pt}{22pt}\selectfont
LaTeX im Vergleich zu klassischen Textverarbeitungen
}

\vspace{0.4cm}

Schwerpunkt: Typografie, Dokumentstruktur und reproduzierbare Dokumenterstellung

\vspace{0.8cm}

\section*{Einordnung}

Dieses Dokument vergleicht LaTeX mit klassischen Textverarbeitungen im Kontext strukturierter technischer Dokumente.

Der Fokus liegt weniger auf einer allgemeinen Bewertung einzelner Werkzeuge, sondern auf den unterschiedlichen Ansätzen bei:

\begin{itemize}
\item Typografie
\item Strukturierung
\item Reproduzierbarkeit
\item Versionskontrolle
\item Wartbarkeit
\item Automatisierung
\end{itemize}

\section*{Vergleich ausgewählter Eigenschaften}

\renewcommand{\arraystretch}{1.35}

\begin{longtable}{>{\raggedright\arraybackslash}p{5.5cm} >{\raggedright\arraybackslash}p{4.3cm} >{\raggedright\arraybackslash}p{4.3cm}}

\toprule
\textbf{Eigenschaft} & \textbf{LaTeX} & \textbf{Klassische Textverarbeitung} \\
\midrule
\endhead

Typografische Konsistenz & \good{Sehr hoch} & \neutralwarn{Abhängig von Vorlage und Arbeitsweise} \\

Trennung von Inhalt und Layout & \good{Konsequent strukturiert} & \neutralwarn{Oft vermischt} \\

Versionierung mit Git & \good{Sehr gut geeignet} & \neutralwarn{Eingeschränkt praktikabel} \\

Reproduzierbare Ausgabe & \good{Deterministische PDF Erzeugung} & \neutralwarn{Teilweise systemabhängig} \\

Automatisierung & \good{Sehr gut integrierbar} & \neutralwarn{Teilweise möglich} \\

Große Dokumente & \good{Sehr stabil} & \neutralwarn{Abhängig von Struktur und Umfang} \\

Mikrotypografie & \good{Direkt integrierbar} & \neutralwarn{Teilweise verfügbar} \\

Semantische Strukturierung & \good{Zentraler Bestandteil} & \neutralwarn{Oft visuell statt strukturell} \\

Wiederverwendbarkeit von Inhalten & \good{Modular möglich} & \neutralwarn{Eher manuell} \\

Layoutkontrolle & \good{Sehr präzise} & \neutralwarn{Visuell orientiert} \\

Einarbeitungszeit & \neutralwarn{Höher} & \good{Niedrige Einstiegshürde} \\

Spontane visuelle Bearbeitung & \neutralwarn{Weniger direkt} & \good{Sehr komfortabel} \\

\bottomrule

\end{longtable}

\section*{Typografische Beispiele}

\section*{Dokumentenstruktur und Wartbarkeit}

Ein wesentlicher Unterschied liegt nicht nur im finalen Erscheinungsbild, sondern in der Struktur des Dokumentes.

LaTeX eignet sich besonders für:

\begin{itemize}
\item reproduzierbare Dokumentgenerierung
\item Versionskontrolle mit Git
\item modulare Dokumentstrukturen
\item automatisierte PDF Builds
\item langfristig wartbare Dokumente
\item Trennung von Inhalt und Gestaltung
\end{itemize}

Dadurch entsteht ein Workflow, der eher mit Softwareentwicklung oder technischer Dokumentation vergleichbar ist als mit klassischer visueller Textbearbeitung.

\section*{Fazit}

LaTeX und klassische Textverarbeitungen verfolgen unterschiedliche Ansätze.

Klassische Textverarbeitungen bieten Vorteile bei schneller visueller Bearbeitung und geringer Einstiegshürde.

LaTeX zeigt seine Stärken besonders bei:

\begin{itemize}
\item typografischer Konsistenz
\item strukturierten Dokumenten
\item reproduzierbaren Ergebnissen
\item Automatisierung
\item und langfristiger Wartbarkeit
\end{itemize}

Die Wahl des geeigneten Werkzeugs hängt letztlich vom jeweiligen Einsatzzweck, der Dokumentkomplexität und dem gewünschten Workflow ab.

\end{document}
```

Wichtige Änderungen gegenüber der ursprünglichen Version:

* Neutralere Sprache statt absoluter Bewertungen.
* Weniger konfrontative Farbgebung.
* Fokus auf Struktur und Reproduzierbarkeit statt nur auf „bessere Typografie“.
* Ergänzung von Git, Automatisierung und Wartbarkeit.
* Professionellere Bezeichnung „Klassische Textverarbeitung“.
* Stärkere Einordnung als technischer Workflow statt Werkzeugvergleich.
* Differenzierteres Fazit ohne Gewinner Verlierer Darstellung.
